\documentclass[11pt,a4paper]{article}
\usepackage{oi-statement}
\OISetup{CSP-J 模拟赛}{第六套}{入门级}
\OISetCompileOptions{-O2 -std=c++14}
\begin{document}
\OICoverTitle
\OICoverTable{4}{
题目名称 & 最多求余 & 各乘一个 & 三元组 & 切割\\\hline

题目类型 & 传统型 & 传统型 & 传统型 & 传统型\\\hline

目录 & {\ttfamily maxmod} & {\ttfamily crossmul} & {\ttfamily triples} & {\ttfamily splitnum}\\\hline

可执行文件名 & {\ttfamily maxmod} & {\ttfamily crossmul} & {\ttfamily triples} & {\ttfamily splitnum}\\\hline

输入文件名 & {\ttfamily maxmod.in} & {\ttfamily crossmul.in} & {\ttfamily triples.in} & {\ttfamily splitnum.in}\\\hline

输出文件名 & {\ttfamily maxmod.out} & {\ttfamily crossmul.out} & {\ttfamily triples.out} & {\ttfamily splitnum.out}\\\hline

每个测试点时限 & 1.0 秒 & 1.0 秒 & 1.0 秒 & 1.0 秒\\\hline

内存限制 & 512 MiB & 512 MiB & 512 MiB & 512 MiB\\\hline

测试点数目 & 10 & 10 & 10 & 10\\\hline

测试点是否等分 & 是，每点 10 分 & 是，每点 10 分 & 是，每点 10 分 & 是，每点 10 分\\\hline
}
\OICompilation{4}{
对于 C++ 语言 & {\ttfamily maxmod.cpp} & {\ttfamily crossmul.cpp} & {\ttfamily triples.cpp} & {\ttfamily splitnum.cpp}\\\hline
}
\OICoverNotes
\begin{enumerate}
\item 本套为 CSP-J 训练模拟赛，共 4 题，每题 100 分，总分 400 分。
\item 各题均采用标准输入输出，提交对应的同名 C++ 源程序。
\item 每题有 10 个测试点，每点 10 分，按测试点独立计分。
\item 答案采用普通文本比较，忽略行末空格和文件末尾多余换行。
\item 题面样例与下发大样例不计入正式测试点。
\item 程序可使用的栈空间与题目的内存限制一致。
\item 样例框左侧行号不属于输入输出。
\item 使用 C++14，编译选项为 \texttt{-O2 -std=c++14}。
\end{enumerate}

\OIProblem{最多求余}{maxmod}

\OISection{题目描述}

给定 $n$ 个正整数 $a_1,a_2,\ldots,a_n$。将每个数对 $k$ 取余，求出现次数最多的余数。

如果有多个余数并列出现最多，输出其中最小的余数。余数可以为 $0$。

\OISection{输入格式}

第一行两个整数 $n,k$。第二行 $n$ 个整数 $a_1,\ldots,a_n$。

\OISection{输出格式}

输出一个整数，表示所求余数。

\begin{OISample}{1}{输入}
4 3
7 10 2 16
\end{OISample}
\begin{OISample}{1}{输出}
1
\end{OISample}
\OISection{样例 1 解释}

余数依次为 $1,1,2,1$，所以答案是 $1$。

\begin{OISample}{2}{输入}
2 10
12 13
\end{OISample}
\begin{OISample}{2}{输出}
2
\end{OISample}
\OISection{样例 2 解释}

余数为 $2,3$，次数相同，选择较小的 $2$。

\OIAdditionalSample{3}{maxmod/maxmod}
\OIAdditionalSample{4}{maxmod/maxmod}
\Needspace{7\baselineskip}
\OISection{数据范围}

$1\le n\le10^4$，$1\le k\le10^5$，$1\le a_i\le10^7$。前两组条件按累计比例分别占 20\%、60\%。

每个测试点 10 分，共 10 个测试点，总分 100 分。

\begin{OIDataTable}{|c|c|c|}
测试点 & 本组分值 & 额外限制\\\hline
1～2 & 20 & 所有元素对 $k$ 的余数相同\\\hline
3～6 & 40 & 所有实际出现的余数，其出现次数两两不同\\\hline
7～10 & 40 & 无额外限制\\\hline
\end{OIDataTable}

\OIProblem{各乘一个}{crossmul}

\OISection{题目描述}

给定长度为 $n$ 的整数数组 $a$ 和长度为 $x$ 的整数数组 $b$。

从 $a$ 中恰好选择 $x$ 个位置 $i_1<i_2<\cdots<i_x$，得到分数

\[
\sum_{j=1}^{x} a_{i_j}b_j.
\]

求分数的最大值。选中的位置不必连续，但不能重复使用一个位置，也不能改变 $b$ 的顺序。数组元素可以是正数、负数或零。

\OISection{输入格式}

第一行两个整数 $n,x$。第二行 $n$ 个整数 $a_i$。第三行 $x$ 个整数 $b_j$。

\OISection{输出格式}

输出一个整数，表示最大分数。

\begin{OISample}{1}{输入}
5 2
1 5 2 4 3
2 -1
\end{OISample}
\begin{OISample}{1}{输出}
8
\end{OISample}
\OISection{样例 1 解释}

选择 $a_2=5$ 和 $a_3=2$，分数为 $5\times2+2\times(-1)=8$。

\begin{OISample}{2}{输入}
3 3
-2 -5 -1
1 2 3
\end{OISample}
\begin{OISample}{2}{输出}
-15
\end{OISample}
\OISection{样例 2 解释}

必须选满三个数，所以分数为 $(-2)\times1+(-5)\times2+(-1)\times3=-15$。

\OIAdditionalSample{3}{crossmul/crossmul}
\OIAdditionalSample{4}{crossmul/crossmul}
\Needspace{9\baselineskip}
\OISection{数据范围}

$1\le n\le10^5$，$1\le x\le\min(n,10)$，$-10^8\le a_i,b_j\le10^8$。各行独立计分，不是累计包含关系。

每个测试点 10 分，共 10 个测试点，总分 100 分。

\begin{OIDataTable}{|c|c|c|}
测试点 & 本组分值 & 额外限制\\\hline
1～2 & 20 & $n\le20$\\\hline
3～4 & 20 & $x=1$\\\hline
5～6 & 20 & $x=2$\\\hline
7～8 & 20 & $x=3$\\\hline
9～10 & 20 & 无额外限制\\\hline
\end{OIDataTable}

\OIProblem{三元组}{triples}

\OISection{题目描述}

有 $n$ 个三元组，第 $i$ 个为 $(a_i,b_i,c_i)$。

对于下标对 $(i,j)$，其中 $1\le i\le j\le n$，如果

\[
2\min(a_i+a_j,b_i+b_j)\le\max(a_i+a_j,b_i+b_j),
\]

则称它为一个好的三元组对，产生价值 $c_ic_j$。

求所有好的三元组对的总价值对 $10^9+7$ 取模的结果。允许 $i=j$；同一下标对只统计一次。

\OISection{输入格式}

第一行一个整数 $n$。接下来 $n$ 行，每行三个整数 $a_i,b_i,c_i$。

\OISection{输出格式}

输出总价值对 $10^9+7$ 取模的结果。

\begin{OISample}{1}{输入}
3
2 6 5
3 4 1
5 6 2
\end{OISample}
\begin{OISample}{1}{输出}
30
\end{OISample}
\OISection{样例 1 解释}

只有 $(1,1)$ 和 $(1,2)$ 满足条件，总价值为 $5\times5+5\times1=30$。

\OIAdditionalSample{2}{triples/triples}
\OIAdditionalSample{3}{triples/triples}
\Needspace{7\baselineskip}
\OISection{数据范围}

$1\le n\le10^5$，$1\le a_i,b_i,c_i\le10^9$。前两档累计分别为 20\%、40\%。

每个测试点 10 分，共 10 个测试点，总分 100 分。

\begin{OIDataTable}{|c|c|c|}
测试点 & 本组分值 & 额外限制\\\hline
1～2 & 20 & $n\le1000$，且取模前的总价值小于 $10^9+7$\\\hline
3～4 & 20 & $n\le1000$\\\hline
5～10 & 60 & 无额外限制\\\hline
\end{OIDataTable}

\OIProblem{切割}{splitnum}

\OISection{题目描述}

给定 $n$ 个正整数。你可以在任意一个原数的十进制表示内部，选择若干个相邻数位间隙切开。每切开一个间隙计一次操作。

例如，\texttt{\detokenize{123456}} 可以切成 \texttt{\detokenize{123}} 和 \texttt{\detokenize{456}}，也可以切成 \texttt{\detokenize{1}}、\texttt{\detokenize{23}} 和 \texttt{\detokenize{456}}。不能改变数位顺序，也不能把两个原数拼接。

每一段按十进制整数解释；允许切出的数段有前导零，例如 \texttt{\detokenize{100}} 切成 \texttt{\detokenize{1}} 和 \texttt{\detokenize{00}} 后，两段的数值为 $1,0$。

求使所有切割后的整数之和为 $k$ 的倍数，最少需要的切割次数。可以不切割；如果无法做到，输出 $-1$。

\OISection{输入格式}

第一行两个正整数 $n,k$。第二行 $n$ 个正整数 $a_i$，原数的输入不含前导零。

\OISection{输出格式}

输出最少切割次数；如果无解，输出 \texttt{\detokenize{-1}}。

\begin{OISample}{1}{输入}
2 5
12 2
\end{OISample}
\begin{OISample}{1}{输出}
1
\end{OISample}
\OISection{样例 1 解释}

将 \texttt{\detokenize{12}} 切成 \texttt{\detokenize{1}} 和 \texttt{\detokenize{2}}，再加另一个 $2$，总和为 $5$，答案为 $1$。

\begin{OISample}{2}{输入}
2 3
123 6
\end{OISample}
\begin{OISample}{2}{输出}
0
\end{OISample}
\OISection{样例 2 解释}

原和为 $129$，已经是 $3$ 的倍数，不需要切割。

\begin{OISample}{3}{输入}
2 2
1 2
\end{OISample}
\begin{OISample}{3}{输出}
-1
\end{OISample}
\OISection{样例 3 解释}

两个数都是一位数，不能再切。原和为 $3$，无法成为 $2$ 的倍数。

\OIAdditionalSample{4}{splitnum/splitnum}
\OIAdditionalSample{5}{splitnum/splitnum}
\Needspace{7\baselineskip}
\OISection{数据范围}

$1\le n,k\le400$，$1\le a_i\le10^{18}$。前 20\% 满足第一行条件；前 60\% 均满足第二行条件，其中前 20\% 同时满足第一行。

每个测试点 10 分，共 10 个测试点，总分 100 分。

\begin{OIDataTable}{|c|c|c|}
测试点 & 本组分值 & 额外限制\\\hline
1～2 & 20 & $n,k\le10$\\\hline
3～6 & 40 & $n,k\le100$，$a_i\le10^5$\\\hline
7～10 & 40 & 无额外限制\\\hline
\end{OIDataTable}

\end{document}
